
\section{11. Conclusion}

\subsection{11.1 Summary}

We presented a production-hardened implementation of **Grouped-Query 
Attention** (GQA) supporting flat, nested, causal, and GQA-native 
architectures. GQA reduces KV cache by $H / G$ at inference with minimal 
quality loss.

\subsection{11.2 Results Recap}

\begin{table}[H]
\centering
\small
\begin{tabular}{llll}
\toprule
\textbf{Model} & \textbf{Acc $\Delta$} & \textbf{KV ratio} & \textbf{Params $\Delta$} \\
\midrule
DistilBERT & -2.2\% & **2.0$\times$** & 1.05$\times$ \\
RoBERTa & -3.1\% & **2.0$\times$** & 1.05$\times$ \\
GPT-2 & -2.4\% & **2.0$\times$** & 1.05$\times$ \\
\bottomrule
\end{tabular}
\end{table}

\subsection{11.3 Key Takeaways}

1. \textbf{GQA reduces KV cache, not parameters.} For parameter reduction, use 
   A3 or AASVD.
2. \textbf{Quality loss is graceful.} At $G = H/2$, drops are small.
3. \textbf{Long-context benefit is linear.} Savings grow with sequence length.
4. \textbf{Compatible with quantization} for multiplicative savings.

\subsection{11.4 Future Work}

\subsubsection{Short-Term}
\begin{itemize}
\item Fused kernels for KV repetition
\item Long-context evaluation ($L = 4096$)
\item Per-head grouping analysis
\item Non-uniform grouping
\end{itemize}

\subsubsection{Long-Term}
\begin{itemize}
\item Theoretical expressiveness analysis
\item Cross-lingual evaluation
\item Integration with MoE
\end{itemize}

\subsection{11.5 Closing Remarks}

GQA is a mature technique with wide adoption in industry (LLaMA, Mistral, 
Qwen). Our implementation makes it accessible to researchers working with 
diverse architectures. We release all code and checkpoints.

\subsection{11.6 Acknowledgments}

We thank {{ACKNOWLEDGMENTS}}.

\subsection{11.7 Author Contributions}

\begin{itemize}
\item \textbf{{{AUTHOR\_1}}} implemented the code, ran experiments, wrote the paper.
\item \textbf{{{AUTHOR\_2}}} provided feedback, edited.
\end{itemize}

\subsection{11.8 Conflict of Interest}

None.